\section{Worked exact finite examples}\label{sec:examples}

The reference programs are not proofs in iPRA.  They serve two narrower
purposes: they test the normal-form algorithms exactly, and they produce small
certificates against which the future object formalization can be regressed.
All arithmetic in the examples is integral or rational; no floating point is
used.

\subsection{A path where greedy choice fails}

Use the three-level relation from \cref{thm:whole-path}.  The level sizes are
\((2,2,2)\); with little-endian base two, the path-code bound is \(2^3=8\).
The transitions are
\[
  0\to0,\quad 1\to1
  \qquad\text{at stage }0,
\]
and
\[
  1\to0
  \qquad\text{at stage }1.
\]
Both initial states are allowed and the final state must be zero.

The greedy selector chooses initial state zero, follows \(0\to0\), and has no
second edge.  Whole-code search checks codes \(0,\ldots,7\) and selects code
\(3\).  Its digits are
\[
  \operatorname{dig}_2(3,0)=1,
  \quad
  \operatorname{dig}_2(3,1)=1,
  \quad
  \operatorname{dig}_2(3,2)=0.
\]
Thus the selected path is \((1,1,0)\), which satisfies every local and terminal
condition.  The complete machine-readable run is
\pathref{examples/whole_path_certificate.json}.

This example validates the framework's insistence that local least choice and
complete finite choice are different contracts.

\subsection{Exact rational bisection of \texorpdfstring{\(y^3+y=x\)}{y^3+y=x}}

Let
\[
  f_x(y)=y^3+y-x.
\]
For rationals \(u>v\),
\[
 f_x(u)-f_x(v)
 =(u-v)(u^2+uv+v^2+1)>0,
\]
so \(f_x\) is strictly increasing.  A sign bracket
\(f_x(\ell)\le0\le f_x(u)\) therefore contains the unique real solution.

At each step let \(m=(\ell+u)/2\).  If \(f_x(m)\le0\), replace \(\ell\) by
\(m\); otherwise replace \(u\) by \(m\).  The sign bracket is preserved, and
the width is halved.

\begin{theorem}[Bisection contract]
After \(p\) exact rational steps from a sign bracket \([\ell_0,u_0]\), the
resulting bracket \([\ell_p,u_p]\) satisfies
\[
 f_x(\ell_p)\le0\le f_x(u_p),
 \qquad
 u_p-\ell_p=\frac{u_0-\ell_0}{2^p}.
\]
\end{theorem}

\begin{proof}
Induct on \(p\).  The base is the supplied bracket.  At the successor step,
the exact comparison with zero selects one half.  Monotonicity gives the new
sign inequalities, and elementary rational arithmetic gives half the previous
width.  The invariant consists entirely of polynomial equalities,
inequalities, denominator positivity, and a width equation, so it has a QF
coding in iPRA.
\end{proof}

The included run takes \(x=2\), initial interval \([0,2]\), and twelve steps.
The exact result is
\[
  [\ell_{12},u_{12}]
   =\left[1,\frac{2049}{2048}\right],
  \qquad
  u_{12}-\ell_{12}=\frac1{2048}.
\]
The first midpoint is exactly \(1\), where \(1^3+1=2\); the implementation
retains equality on the lower side, so later intervals continue to contain the
root at their left endpoint.  Every local transition and the terminal width
are independently rechecked in
\pathref{examples/cubic_bisection_certificate.json}.

\subsection{FQNF encoding of the bisection trace}

There are two natural certificates.

\begin{enumerate}
\item Store the \(p\) branch bits in a code below \(2^p\).  A verifier
reconstructs the rational endpoints by primitive recursion and checks the sign
comparison chosen at every step.
\item Store all endpoint numerators and denominators in a mixed-radix master
code under explicit growth bounds.  The verifier checks every transition
locally.
\end{enumerate}

The branch-bit certificate is smaller because the trace is deterministic once
\(x,\ell_0,u_0\) are fixed.  The full-state trace is a useful generic format
when a producer is external or the transition is relational.

\subsection{Regression tests}

The exact test suite checks:

\begin{itemize}
\item exhaustive uniform-base round trips for bases two through six and lengths
up to four;
\item exhaustive mixed-radix round trips for radices \((2,3,4,2)\);
\item least-certificate selection and failure when no certificate exists;
\item complete bounded-list enumeration;
\item greedy path failure and whole-code success;
\item stability of independently chosen rational intervals;
\item every local cubic-bisection transition and the exact final width;
\item backward demand through a fixed pipeline.
\end{itemize}

All eleven tests pass in the delivered artifact.  They are regression evidence
for the finite algorithms, not substitutes for Lean kernel checking.
